\section{Experimental Setup}

\subsection{Datasets and Data Preprocessing}

\subsubsection{Sub-Problem 1 Vision Dataset}
The camera vehicle estimation dataset consists of \textbf{6,000 manually annotated traffic images} captured across surveillance cameras from the Ho Chi Minh City Traffic Portal. The dataset is partitioned into:
\begin{itemize}
    \item \textbf{Training Set:} 4,000 images.
    \item \textbf{Validation Set:} 1,000 images.
    \item \textbf{Testing Set:} 1,000 images.
\end{itemize}
Ground-truth annotations provide individual continuous counts for passenger cars and motorbikes. Images are resized to $224 \times 224$ pixels and normalized with ImageNet mean and standard deviation.

\subsubsection{Sub-Problem 2 Traffic Graph Time Series Dataset}
Using Stage 1 predictions across 608 active camera nodes, time series are aggregated at 5-minute intervals ($step\_minutes=5$). Missing steps ($\le 30$ minutes) are linearly interpolated and smoothed via double-rolling transformation ($window_1=3$, $window_2=5$).

Historical lookback length is $T_{\text{in}} = 24$ steps (120 minutes history), and forecast horizon is $H = 6$ steps (30 minutes ahead).

The time-series dataset is split chronologically:
\begin{itemize}
    \item \textbf{Training Set:} First 80\% of the time range.
    \item \textbf{Validation Set:} Middle 10\% of the time range.
    \item \textbf{Testing Set:} Final 10\% of the time range.
\end{itemize}

\subsection{Hyperparameter Settings and Training Details}

All models across both sub-problems are implemented in PyTorch 2.x and trained hardware-accelerated on an **NVIDIA T4 GPU (16 GB VRAM)**.

\subsubsection{Sub-Problem 1 Training Hyperparameters}
The vision backbones (ConvNeXt-Tiny \cite{convnext} and ResNet-50 \cite{resnet}) with custom MLP regression heads are trained under the following setup:
\begin{itemize}
    \item \textbf{Hardware Accelerator:} NVIDIA T4 GPU (16 GB VRAM).
    \item \textbf{Optimizer:} AdamW optimizer with initial learning rate $\eta = 1 \times 10^{-4}$ and weight decay coefficient $1 \times 10^{-2}$.
    \item \textbf{Learning Rate Schedule:} Cosine Annealing Learning Rate Scheduler decaying down to $\eta_{\min} = 1 \times 10^{-6}$.
    \item \textbf{Epochs and Batch Size:} Models are trained for 100 epochs with a mini-batch size of $B = 32$.
    \item \textbf{Loss Function:} Smooth $L_1$ (Huber) Loss \cite{huber} with threshold $\beta = 1.0$.
    \item \textbf{Data Augmentations:} Random horizontal flips (probability $p=0.5$), color jittering ($\pm 15\%$ brightness, contrast, and saturation), and random spatial cropping resized to $224 \times 224$.
\end{itemize}

\subsubsection{Sub-Problem 2 Training Hyperparameters}
The spatio-temporal graph models (GCN-LSTM \cite{gcn, tgcn}, STGCN Baseline \cite{stgcn}, and TA-STGCN \cite{attention}) are trained under the following setup:
\begin{itemize}
    \item \textbf{Hardware Accelerator:} NVIDIA T4 GPU (16 GB VRAM).
    \item \textbf{Optimizer:} Adam optimizer \cite{adam} with initial learning rate $\eta = 5 \times 10^{-4}$ and weight decay coefficient $1 \times 10^{-4}$.
    \item \textbf{Learning Rate Schedule:} ReduceLROnPlateau scheduler decaying the learning rate by a factor of $0.5$ if validation loss fails to improve for 10 consecutive epochs ($\text{patience}=10$).
    \item \textbf{Max Epochs and Early Stopping:} Trained for up to 500 epochs with an early stopping patience of 30 epochs based on Validation MAE.
    \item \textbf{Batch Size:} Training mini-batch size $B = 16$ (evaluation batch size $\min(B, 32) = 16$).
    \item \textbf{Loss Function:} Pure Huber Loss \cite{huber} ($\delta = 1.0$) trained end-to-end.
    \item \textbf{Multi-Seed Evaluation Protocol:} Models are evaluated across \textbf{5 independent random seeds}: $\mathcal{S} = \{42, 100, 2024, 777, 999\}$. All reported metrics present $\text{Mean} \pm \text{Std}$ over the 5 seeds.
\end{itemize}

\subsection{Evaluation Metrics}

Model performance is evaluated using Mean Absolute Error (MAE), Root Mean Squared Error (RMSE), and Mean Squared Error (MSE):
\begin{equation}
\text{MAE} = \frac{1}{N \cdot H \cdot M} \sum_{i=1}^{N} \sum_{h=1}^{H} \sum_{k=1}^{M} \left| y_{i,h}^{(k)} - \hat{y}_{i,h}^{(k)} \right|,
\end{equation}
\begin{equation}
\text{RMSE} = \sqrt{ \frac{1}{N \cdot H \cdot M} \sum_{i=1}^{N} \sum_{h=1}^{H} \sum_{k=1}^{M} \left( y_{i,h}^{(k)} - \hat{y}_{i,h}^{(k)} \right)^2 },
\end{equation}
\begin{equation}
\text{MSE} = \frac{1}{N \cdot H \cdot M} \sum_{i=1}^{N} \sum_{h=1}^{H} \sum_{k=1}^{M} \left( y_{i,h}^{(k)} - \hat{y}_{i,h}^{(k)} \right)^2,
\end{equation}
where $N=608$ graph nodes, $H=6$ forecast steps, and $M$ is the test batch count.
